\documentclass[12pt,a4paper]{article}
\usepackage[utf8]{inputenc}
\usepackage{geometry}
\usepackage{fancyhdr}
\usepackage{hyperref}

% Geometría de la página
\geometry{
 a4paper,
 total={170mm,257mm},
 left=20mm,
 top=20mm,
}

% Encabezado y pie de página
\pagestyle{fancy}
\fancyhf{}
\fancyhead[C]{Propuesta de Webinar: Seguridad en Criptoactivos}
\fancyfoot[C]{\thepage}
\renewcommand{\headrulewidth}{0.4pt}
\renewcommand{\footrulewidth}{0.4pt}

% Título
\title{
    \vspace{2cm}
    \textbf{Propuesta de Webinar Gratuito} \\
    \large \textit{Seguridad en Criptoactivos: USDT y Monedas Estables como Herramienta Financiera}
    \vspace{1.5cm}
}
\author{Ing. Daniel Velarde Kubber}
\date{La Paz, 13 de junio de 2025}

\begin{document}

\maketitle
\thispagestyle{empty}
\clearpage

\section*{1. Detalle General de la Propuesta}

\begin{itemize}
    \item \textbf{Título del Webinar:} Seguridad en Criptoactivos: USDT y Monedas Estables como Herramienta Financiera.
    \item \textbf{Propuesta dirigida a:} Universidad Real de la Cámara Nacional de Comercio.
    \item \textbf{Ponente:} Ing. Daniel Velarde Kubber - Ingeniero Industrial y de Sistemas, entusiasta en Data y Blockchain.
    \item \textbf{Fecha propuesta:} 13 de junio de 2025.
    \item \textbf{Ciudad:} La Paz, Bolivia.
    \item \textbf{Modalidad:} Virtual (a través de la plataforma que disponga la universidad).
    \item \textbf{Duración Total:} 1 hora (45 minutos de exposición y 15 minutos para preguntas y respuestas).
    \item \textbf{Costo:} Gratuito.
\end{itemize}

\section*{2. Objetivo del Webinar}

El objetivo principal es proporcionar a la comunidad universitaria una comprensión clara y precisa sobre las monedas estables (stablecoins), con un enfoque especial en Tether (USDT). Se busca educar a los participantes sobre su funcionamiento, su utilidad en el contexto económico actual de Bolivia, y, fundamentalmente, sobre las mejores prácticas de seguridad para mitigar los riesgos inherentes a su uso y promover una adopción responsable.

\section*{3. Temario y Contenido (Duración: 45 minutos)}

\begin{enumerate}
    \item \textbf{Introducción a los Criptoactivos (5 min)}
    \begin{itemize}
        \item Breve panorama del ecosistema cripto.
        \item Diferencia entre criptomonedas volátiles y monedas estables.
    \end{itemize}
    
    \item \textbf{¿Qué es una Moneda Estable (Stablecoin)? (10 min)}
    \begin{itemize}
        \item Definición y propósito: Anclaje al valor de activos del mundo real.
        \item Tipos de stablecoins: Colateralizadas en fiat (USDT, USDC), en cripto (DAI) y algorítmicas.
        \item \textbf{Foco principal: ¿Qué es USDT (Tether)?} Funcionamiento y respaldo.
    \end{itemize}
    
    \item \textbf{Relevancia en la Situación Actual de Bolivia (10 min)}
    \begin{itemize}
        \item El uso de USDT como alternativa para la preservación de valor frente a la escasez de divisas.
        \item Facilidad para transacciones internacionales y remesas.
        \item Ventajas y desafíos en el contexto local.
    \end{itemize}
    
    \item \textbf{Seguridad y Riesgos Latentes (15 min)}
    \begin{itemize}
        \item \textbf{Riesgos clave:} Volatilidad del colateral, centralización, regulatorios y de contraparte.
        \item \textbf{Medidas de Seguridad Esenciales para el Usuario:}
        \begin{itemize}
            \item Elección de billeteras (wallets) seguras: calientes vs. frías.
            \item Custodia personal: La importancia de las claves privadas.
            \item Identificación de estafas comunes (phishing, esquemas Ponzi).
            \item Uso de plataformas de intercambio (exchanges) reconocidas y seguras.
        \end{itemize}
        \item La importancia de la debida diligencia y el uso responsable.
    \end{itemize}
\end{enumerate}

\section*{4. Sesión de Preguntas y Respuestas (15 minutos)}

Al finalizar la presentación, se abrirá un espacio para que los asistentes puedan realizar preguntas y resolver sus dudas directamente con el ponente.

\vspace{2cm}

\begin{center}
    Agradezco de antemano su tiempo y consideración. \\
    \vspace{0.5cm}
    \textbf{Ing. Daniel Velarde Kubber} \\
    \textit{Entusiasta en Data y Blockchain}
\end{center}

\end{document}